\scanpage{55}
\mspage{7.}
nent de $D$ mod $\mathbb C^*$, donc qui proviennent de $D\times\mathbb C^*$. Ceux qui conservent \deleted{\unsure{P}} chaque $P_\omega$, i.e. chaque sous-espace isotrope $V_{\mathbb C}^\omega$, sont ceux qui proviennent de $D^+\times\mathbb C^*$.

Si on plonge
\[
(36)\qquad D^+\times\mathbb C^*\unsure{\lhook\joinrel\relbar|\hookrightarrow}(\mathbb C^*)^{\underline\omega}
\]
comme dans (34), alors l'image est formée des éléments \inserted{$(\lambda_\omega)_{\omega\in\underline\omega}$} tels que, choisissant $\omega\in\underline\omega$, on ait
\[
(37)\qquad\deleted{\unsure{\prod\lambda_\omega}}(\lambda_\omega\lambda_{\omega'}^{-1})^n=1
\]
(condition indépendante du choix de $\omega$), ce qui, pour \deleted{$\lambda_\omega,\lambda_{\omega'}$} $\lambda_\omega,\lambda_{\omega'}\in\mu_m=\mu_{2n}$, équivaut à la condition sous la forme plus symétrique \unsure{en apparence} (35). Donc les éléments de $\mu_m^{\underline\omega}$ qui ne satisfont pas cette relation ne stabilisent pas $V^{**}$.

On s'intéresse à des choix de pts \unsure{base} qui, en \deleted{tout} état de cause, soient stables par $D\times\{1,\tau\}$, et éventuellement $=$ par le groupe de symétries plus grand
\[
(D\times\mu_m)/\{\pm1\}\times\{1,\tau\}
\]
(d'ordre $72$ dans le cas $n=3$)
\mspage{54}
